\section{Construction in characteristic zero}\label{con:section}

We now produce the enhanced eigencomplex required by
Proposition~\ref{desc:ordinary}.  Starting with a parameter on a complex
curve punctured at $s$ points, we attach one second component at all
$s$ punctures and smooth the resulting nodes.  The unramified
correspondence supplies an eigencomplex on the smooth generic curve.
Its nearby cycles produce the enhanced flags, and torus descent gives
the ordinary-flag eigensheaf.  We follow the construction of
\cite[Section~8]{CT}, explaining the compatibility at several nodes and
the nonvanishing on the specified component of the bundle stack.

\subsection{The unramified input}

We use the following precise input from the characteristic-zero
correspondence.  Its local constructibility conclusion is essential
for applying nearby cycles.

\begin{lemma}[{\cite[Lemma~8.1]{CT}}]\label{con:unramified}
Let $C$ be a smooth projective connected curve over an algebraically
closed field of characteristic zero, and let $G$ be simple and simply
connected.  A continuous $\widehat G$-local system $\tau$ on $C$,
defined over a finite extension of $\mathbb Q_\ell$ and with
Zariski-dense image, admits a nonzero locally bounded constructible
geometric $E$-adic eigencomplex on $\Bun_G(C)$.  Its Hecke
eigenstructure has the relative Satake normalization and is coherent
for all finite sets of legs, including collision maps.
\end{lemma}

We recall why the input includes these finiteness and coherence
assertions; we use the full correspondence only through this lemma.
The characteristic-zero theorem of Gaitsgory--Raskin identifies the
automorphic nilpotent category with the ind-coherent category with
nilpotent singular support on the restricted stack of
$\widehat G$-local systems, compatibly with its quasi-coherent spectral
action \cite[Main Theorem~1.3.9(ii), Lemma~1.3.7, and
Theorem~1.1.7]{GR}.  At a dense parameter the adjoint group
$\widehat G$ has trivial centralizer.  Consequently
$H^0(C,\ad\tau)=0$, and duality gives $H^2(C,\ad\tau)=0$.
The corresponding formal component is formally smooth, with no automorphisms;
its ind-coherent skyscraper is nonzero, compact, and has zero
obstruction-theoretic singular support
\cite[Proposition~3.7.2, Corollary~3.7.5, and
Sections~21.1--21.2]{AGKRRV}.  Compactness is preserved by the
correspondence and by inclusion of the automorphic nilpotent category
\cite[Theorem~1.1.7]{GR},
and implies bounded constructibility on every finite-type smooth
chart \cite[Appendix~F, Section~F.2.2]{AGKRRV}.  Finally, the projection
formula for the skyscraper evaluates every universal tensor local
system at $\tau$.  Spectral linearity turns these identities into the
entire Hecke eigenstructure, with its unit, tensor, permutation, and
collision compatibilities
\cite[Main Theorem~14.3.2 and Section~14.3.3]{AGKRRV}.
The Satake labeling is fixed as in Section~\ref{found:section}; the
spectral point is labeled so that evaluation of $V$ is $V_\tau$.
This is the compact-object argument of \cite[Lemma~8.1]{CT}.

\subsection{Several nodes and their bundle types}

For the rest of this section $G=\mathrm{SL}_n$.  Let $X_1/\mathbb C$
be a smooth projective connected curve of genus $g\geq2$, with
$s\geq2$ distinct points $x_{11},\ldots,x_{1s}$.  Write
$D_1=\sum_i x_{1i}$ and $U_1=X_1\setminus|D_1|$, and fix a
continuous dense parameter $\sigma$ on $U_1$, defined over a finite
coefficient field and unipotent at each puncture.  Let
$(X_2,(x_{2i}))$ be a second copy of the marked curve; put
$D_2=\sum_i x_{2i}$ and $U_2=X_2\setminus|D_2|$.  Identify
$x_{1i}$ with $x_{2i}$ for every $i$.  The resulting connected nodal
curve $C_0$ has two components and $s$ nodes.  There is a projective
smoothing
\[
 C\longrightarrow\Spec\mathbb C[[z]]
\]
with smooth connected geometric generic fiber and local equations
$ab=z$ at all nodes.  Indeed, deformations of a nodal curve are
unobstructed, the node-smoothing parameters can be prescribed
independently, and an ample line bundle lifts and algebraizes the
formal deformation.  Taking every parameter to have order one gives
a regular semistable model.

Choose compatible roots $z_m^m=z$, and put $R_m=\mathbb C[[z_m]]$.
The balanced twisted model $C(m)/R_m$ has, at each node, the chart
\begin{equation}\label{con:node-chart}
 \left[\Spec R_m[u,v]/(uv-z_m)\big/\mu_m\right],
 \qquad a=u^m,\quad b=v^m,\quad
 \zeta(u,v)=(\zeta u,\zeta^{-1}v).
\end{equation}
Here the positive action is on the $X_1$ branch.  These charts are
understood \emph{\'etale locally}.  In the regular semistable model the
two components are transverse Cartier divisors; their local equations,
after absorbing a unit, give \'etale maps to $ab=z$ near each node.
Choose neighborhoods containing only their designated node, insert
\eqref{con:node-chart}, and glue with the ordinary curve off the
nodes.  The charts coincide there with their coarse spaces, so the
construction introduces no further stacky points.  It is the
balanced twisted-curve construction of
\cite[Theorems~1.9--1.10, Remark~1.11, and
Proposition~2.2(ii)]{OlssonTwisted}, used in \cite[Section~8.2]{CT}.
For $m\mid m'$ the power maps $u_m=u_{m'}^{m'/m}$ and
$v_m=v_{m'}^{m'/m}$ give compatible transition maps.  Each $C(m)$ is
proper and flat, tame, and has finite diagonal.  Its special-fiber
normalization consists of $X_1$ and $X_2$ with an $m$th root-stack
point at each marking.

Fix a strictly dominant cocharacter $\lambda\in X_*(T)$ and an index
$m$ such that
\begin{equation}\label{con:alcove}
 0<\langle\alpha,\lambda\rangle<m
 \qquad\text{for every positive root }\alpha.
\end{equation}
At every node prescribe the type $\lambda|_{\mu_m}$, using the common
node generator in \eqref{con:node-chart}.  Its centralizer is $T$:
the eigenvalues in the standard representation are pairwise distinct
by \eqref{con:alcove}.  Let $\cB_m$ be the open substack of
$\Bun_G(C(m)/R_m)$ obtained by deleting the other types in the closed
fiber.  There are finitely many types, and their weight multiplicities
are locally constant in families on each residual gerbe, so this is
an open substack.

The stack $\cB_m$ is smooth and locally of finite type over $R_m$.
Algebraicity and local finite presentation follow from the Hom-stack
theorem for a proper flat finitely presented source and target $BG$
with affine diagonal \cite[Theorem~1.2]{HallRydh}.  The obstruction
group at a bundle $P$ is $H^2(C(m)_t,\ad P)$, which vanishes because
coarse pushforward on a tame curve is exact and the coarse space is
one-dimensional.  Its geometric generic fiber is the unlevelled
bundle stack of the smooth generic curve.

Let $B^-$ be the opposite Borel containing $T$.  Define
\[
 \cA_1^+=\Bun_{G,N_B,D_1}(X_1),\qquad
 \cA_2^+=\Bun_{G,R_u(B^-),D_2}(X_2).
\]
Thus the second component has opposite enhanced flags.  The following
description is the several-node form of \cite[Lemma~8.2]{CT}; the
relation between twisted nodes and parahoric level also appears in
\cite[Section~2.3, Proposition~3.1.5, and Lemma~3.2.3]{NYGluing}.

\begin{lemma}\label{con:quotient}
There is an equivalence
\begin{equation}\label{con:special-quotient}
 (\cB_m)_0\simeq
       [(\cA_1^+\times\cA_2^+)/T^s].
\end{equation}
The $i$th torus changes simultaneously the two frames at the $i$th
gluing gerbe, expressed using its common node generator.  A Hecke
modification in $U_j$ pulls back to the usual modification on the
$j$th factor.
\end{lemma}

\begin{proof}
We first calculate on a normalized root disc of the first component.
Put $t=u^m$.  Locally on a test scheme $S$, choose an equivariant
frame in which inertia acts by $\lambda(\zeta)$.  Frames lift from
the origin through infinitesimal neighborhoods by smoothness and
exactness of $\mu_m$-invariants.  Their changes satisfy
\[
 h(\zeta u)=\lambda(\zeta)h(u)\lambda(\zeta)^{-1},
 \qquad h(0)\in T.
\]
In the invariant punctured-disc frame the same change is
$g(t)=\lambda(u)^{-1}h(u)\lambda(u)$.  For a positive root $\alpha$
and $j=\langle\alpha,\lambda\rangle$, its positive and negative root
series transform as
\begin{equation}\label{con:iwahori-series}
 u^j a_\alpha(u^m)\longmapsto a_\alpha(t),
 \qquad
 u^{m-j}b_\alpha(u^m)\longmapsto t b_\alpha(t).
\end{equation}
Toral series are series in $t$.  These are exactly the series in
the Iwahori group
\[
 I_B(S)=\{g\in G(\mathcal O_S[[t]]):g(0)\in B\}.
\]
For completeness the root calculation identifies the whole group:
$h(0)\in T$ puts $h$ in the formal big cell $R_u(B^-)TR_u(B)$;
the formulas in \eqref{con:iwahori-series} transform its root-ordered
factorization into that of $I_B$.  Conversely reduction of an
Iwahori element belongs to $B$, hence to that big cell, and the
inverse formulas give a regular equivariant frame change.

Gluing to the component off its marked points therefore turns a
bundle of this root type into an ordinary $B$-level bundle.
Evaluation $h\mapsto h(0)$ becomes the torus projection $I_B\to T$.
Its kernel is the inverse image of $N_B$ under $I_B\to B$, so
identifying the gerbe restriction with the fixed type torsor is
precisely the extra datum of an enhanced flag.

On the second branch a positive coordinate generator is
$\eta=\zeta^{-1}$.  Its fiber action is
$\lambda(\eta^{-1})$, so the same calculation uses $-\lambda$ and
gives $B^-$.  In the common node generator both fixed type torsors
have automorphism group $T$.  A bundle on the nodal curve consists
of its bundles on the normalization and an identification at every
pair of residual gerbes.  Choosing frames of both gerbe restrictions
makes that identification automatic; changing the two frames
simultaneously leaves it unchanged.  Applying this construction at
the $s$ disjoint pairs gives \eqref{con:special-quotient}.  Multiple
nodes between the same components impose no additional relation:
nodal descent specifies one independent gluing isomorphism at each
node.  With reversed right-action conventions the same quotient
would have inversion on its second torus factor; here we use
simultaneous gerbe frames.

A modification away from the markings transports the gerbe frames
and their identifications.  The corresponding Hecke diagram
therefore pulls back to the diagram on the indicated factor.
\end{proof}

At a leg in the smooth scheme locus of $C(m)$, all node data lie
outside the modification disc.  Disc gluing gives the usual split
affine-Grassmannian models: the bounded maps are proper, the
exact-position maps to bundle and leg are smooth, and all
modifications preserve the chosen types.  The special leg locus is
$U_1\sqcup U_2$.  Thus the family satisfies the geometric hypotheses
of Theorem~\ref{det:specialization} and
Proposition~\ref{found:hecke-specialization}, including successive
modifications and collisions.  Its special cone and dimension bound
are those of the torus quotient in Section~\ref{geom:section}.

\subsection{A parameter on the smoothing}

We must now extend $\sigma$ across the smoothing.  The second copy
carries one parameter whose boundary monodromies are simultaneously
inverse to those on the first.  This is where using an
orientation-reversing homeomorphism of the entire marked surface is
essential.

Choose such a homeomorphism
$f:U_2^{\mathrm{an}}\to U_1^{\mathrm{an}}$ carrying the punctures in
order.  It induces an isomorphism on profinite fundamental groups
by Riemann existence.  If $\rho_1$ represents $\sigma$, put
$\rho_2=\rho_1\circ\widehat f_*$.  A positive loop $c_{2i}$ goes to
a conjugate of $c_{1i}^{-1}$.  Choose paths to the boundary fibers
so that, in these identifications,
\begin{equation}\label{con:boundary-inversion}
              \rho_2(c_{2i})=\rho_1(c_{1i})^{-1}
              \qquad(1\leq i\leq s).
\end{equation}
The choices are made in the full compact image, before taking finite
quotients.  The homeomorphism need not be algebraic: Riemann
existence algebraizes its finite covers.  Figure~\ref{con:gluing-figure}
records the two distinct inversions, of boundary loops and of the
second branch coordinate, that make gluing possible.

The index $m$ remains fixed for the bundle type and the eventual
nearby-cycle calculation on $\cB_m$.  To construct the parameter we
will instead use divisible indices $m_r$, large enough to kill the
finite inertia images at stage $r$.  One fixed root index need not
kill the entire adic boundary monodromy.

\begin{figure}[tb]
\centering
\begin{tikzpicture}[x=1cm,y=1cm,>=Stealth,font=\small]
 \draw[thick] (-2.5,0) ellipse (1.0 and 1.6);
 \draw[thick] (2.5,0) ellipse (1.0 and 1.6);
 \node at (-2.65,0) {$X_1$};
 \node at (2.65,0) {$X_2$};
 \foreach \yy/\ii in {1/1,-1/s} {
   \fill (-1.72,\yy) circle (1.6pt);
   \fill (1.72,\yy) circle (1.6pt);
   \draw[dashed] (-1.72,\yy)--(1.72,\yy);
   \node[above] at (-1.32,\yy) {$x_{1\ii}$};
   \node[above] at (1.32,\yy) {$x_{2\ii}$};
   \node[fill=white,inner sep=2pt] at (0,\yy) {$B\mu_m$};
 }
 \node at (0,0) {$\vdots$};
 \node[align=center] at (-2.5,-2.1)
      {positive branch $u$\\enhanced $B$-flags};
 \node[align=center] at (2.5,-2.1)
      {negative branch $v$\\enhanced $B^-$-flags};
 \draw[->,thick] (1.5,2.05) to[bend right=12]
       node[above,align=center] {$f$ reverses orientation\\at every boundary}
       (-1.5,2.05);
 \node[align=center] at (0,-3.12)
      {$\zeta:(u,v)\mapsto(\zeta u,\zeta^{-1}v)$\qquad
       $\rho_2(c_{2i})=\rho_1(c_{1i})^{-1}$\\[3pt]
       $\displaystyle (\cB_m)_0\simeq
         [(\cA_1^+\times\cA_2^+)/T^s]$};
\end{tikzpicture}
\caption{Normalization and gluing at all $s$ nodes (schematically).
The ovals are the coarse curves of the normalized root stacks.
The dashed pairs are identified, not extra curve components.
For bundle types the gerbes have fixed order $m$ and each pair
contributes one frame-change torus to the displayed quotient.
For parameter covers at stage $r$, replace $m$ by $m_r$ and
$\rho_j$ by its finite quotient at that stage.  The inverse boundary actions
then agree on the gluing gerbe, because its common generator acts
oppositely on the two branch coordinates.}\label{con:gluing-figure}
\end{figure}

\begin{lemma}\label{con:parameter}
Put $K_1=\overline{\mathbb C((z))}$ and let $C_{K_1}$ be the smooth
geometric generic curve.  There is a continuous unramified parameter
$\rho_{K_1}$ on $C_{K_1}$ with the same compact image as $\rho_1$.
Its lattice reductions specialize on $U_1\sqcup U_2$ to those of
$\rho_1$ and $\rho_2$, with the tensor compatibilities of
Proposition~\ref{found:hecke-specialization}.
\end{lemma}

\begin{proof}
Let $H=\rho_1(\pi_1^{\et}(U_1))\subset\Gd(L')$ be the compact
image, for a finite coefficient field $L'$.  Choose a cofinal
decreasing sequence of open normal subgroups $H_r$ and put
$Q_r=H/H_r$.  One obtains such a sequence from congruence kernels
of an $H$-stable lattice in a faithful representation.  Both $U_j$
carry connected $Q_r$-torsors.  Choose indices $m_r$ divisible by
the fixed $m$ and by the orders of all their boundary images, with
$m_r\mid m_{r+1}$.

By the root-stack description of tame covers
\cite[Proposition~A.10]{LieblichOlsson}, the torsors extend to finite
\'etale torsors on the normalization of $C(m_r)_0$.  At each pair
of gerbes their actions agree: the node identifies a positive
generator on the first branch with a negative generator on the
second, and \eqref{con:boundary-inversion} cancels that inversion.
Reduce the chosen boundary-fiber identifications modulo $H_r$ to
identify these restrictions.  They are compatible at every level of
the tower.  The torsors glue to a finite \'etale $Q_r$-torsor
$Y_{r,0}$ on $C(m_r)_0$.  Locally this is nodal gluing of finite
\'etale algebras with identified fibers: on the completed branch
discs the covers are constant, with the same $\mu_{m_r}$-action at
the origin, so their gluing is constant on the nodal cover and
descends equivariantly.  This also proves compatibility under the
power maps of the twisted models.

Proper henselian invariance of finite \'etale covers lifts $Y_{r,0}$
to a torsor $Y_r$ over $C(m_r)$
\cite[Theorem~4.5]{Rydh}.  Its hypotheses hold because the models
are proper with finite diagonal over complete henselian traits.
Full faithfulness lifts the group actions, transition maps after
pullback to divisible models, and all relations between them.
All generic fibers become $C_{K_1}$.  Choose a base section
specializing in $U_1$ and compatible points of the torsor fibers
above it; finite \'etale covers of this strictly henselian section
are constant.  The tower therefore gives a continuous homomorphism
\begin{equation}\label{con:generic-homomorphism}
 \rho_{K_1}:\pi_1^{\et}(C_{K_1})\longrightarrow
                      \varprojlim_r Q_r=H.
\end{equation}

We check its image.  Each $Y_{r,0}$ is connected, since its two
normalization torsors are connected and are glued along nonempty
fibers.  The total torsor $Y_r$ is proper and flat with reduced
geometric fibers.  Its tame coarse space is flat as well: locally
its coordinate ring is the $\mu_{m_r}$-invariant direct summand of an
$R_{m_r}$-flat ring, and invariant formation commutes with base change.
The connected reduced proper special fiber has
$H^0(\mathcal O)=\mathbb C$.  Semicontinuity, applied to the coarse
space, gives $h^0(Y_{r,K_1},\mathcal O)=1$, so the geometric generic
torsor is connected.  Thus \eqref{con:generic-homomorphism} surjects
onto every $Q_r$.  Its image is compact and therefore closed in $H$;
surjectivity to this cofinal system makes it all of $H$.  In
particular the parameter remains Zariski dense.

For an algebraic representation of $\Gd$, enlarge its finite
coefficient field if necessary and choose an $H$-stable lattice.
Each finite reduction factors through some $Q_r$.  Its lifted
torsor supplies a lisse extension over the smooth leg locus after
the corresponding finite trait extension, with the specified
special reductions on both components.  All twisted models agree
away from the nodes.  The tower thus gives the required tensor
compatibilities; maps between lattices are compared after clearing
denominators.  As allowed in
Proposition~\ref{found:hecke-specialization}, different reductions
may require different finite trait extensions.
\end{proof}

\subsection{A nonzero specialization and descent}

Lemma~\ref{con:unramified} now supplies a nonzero locally bounded
constructible eigencomplex $F$ on $\Bun_G(C_{K_1})$ with parameter
$\rho_{K_1}$.  Take its geometric nearby cycles on the fixed model
$\cB_m$.  Propositions~\ref{found:nearby} and
\ref{found:hecke-specialization}, with the geometry and lattice
extensions just verified, give a locally bounded constructible
eigencomplex $\Psi F$ on $(\cB_m)_0$, for legs on both $U_1$ and
$U_2$.  To use it we must prove nonvanishing on precisely this
prescribed-type fiber.

Choose a point of the product in \eqref{con:special-quotient} and
lift its image to a reference bundle $P_0$ over $R_m$.  This is
possible by smoothness of $\cB_m$: lift a point in a smooth chart
over the henselian trait.  Choose also a section $y$ of the smooth
scheme locus, disjoint from the nodes.  Since $F$ is nonzero, some
$K_1$-valued bundle $P$ has nonzero stalk.  Both $P$ and $(P_0)_{K_1}$
are trivial on the affine curve $C_{K_1}\setminus\{y_{K_1}\}$.
Indeed, a rank-$n$ projective module over a Dedekind domain is the
sum of a free rank-$(n-1)$ module and its determinant; an
$\mathrm{SL}_n$-bundle has trivial determinant, and its frame can be
adjusted to respect the chosen determinant trivialization.

An isomorphism off $y_{K_1}$ exhibits $P$ as a modification of
$(P_0)_{K_1}$ in a finite Schubert bound.  Its point is defined over
a finite extension of the trait's fraction field.  Properness of
the bounded Hecke space over the fixed input $P_0$ and leg $y$
extends this modification over the resulting trait.  Its output
still has the prescribed types, because the modification is away
from all nodes.  Lift this section to a finite-type smooth chart
through its special value.  Its generic stalk is the chosen nonzero
stalk, so the closure of the generic nonzero-stalk locus on this
chart meets the special fiber.  Theorem~\ref{det:specialization}
therefore gives
\begin{equation}\label{con:nonzero-nearby}
                         \Psi F\ne0.
\end{equation}
Here finite extensions cause no change to geometric nearby cycles,
as in Proposition~\ref{found:nearby}.

Pull back $\Psi F$ ordinarily along the smooth surjection
$\cA_1^+\times\cA_2^+\to(\cB_m)_0$.  Its pullback is nonzero.
By local constructibility choose a $\mathbb C$-valued point
$(a_1,a_2)$ with nonzero stalk and restrict to
$\cA_1^+\times\{a_2\}$.  This gives a nonzero locally bounded
constructible complex $F_1^+$ on $\cA_1^+$.  By
Lemma~\ref{con:quotient}, the first component's Hecke correspondences
are the base changes of those downstairs and act on the first
factor alone.  Proper base change on each bound makes both
pullbacks compatible with the full coherent Hecke system.  Thus
$F_1^+$ has eigenvalue $\sigma$.  These are ordinary pullbacks;
perversity will be obtained after descent.

\begin{proposition}\label{con:charzero}
Let $X/\mathbb C$ be a smooth projective connected curve of genus
at least two, with $s\geq2$ distinct marked points, divisor $D$ their
sum, and $G=\mathrm{SL}_n$.
Let $\sigma$ be a continuous dense $\Gd$-parameter on their
complement, defined over a finite coefficient extension, with
unipotent monodromy at every marking.  Then
$\Bun_{G,B,D}(X)$ carries a nonzero locally constructible perverse
geometric $E$-adic eigensheaf with eigenvalue $\sigma$, nilpotent
singular support, and the full coherent Hecke system in relative
Satake normalization.
\end{proposition}

\begin{proof}
The construction gives $F_1^+$ on the enhanced stack for $X$.
Proposition~\ref{desc:ordinary} applies because the parameter is
dense and each boundary monodromy is unipotent, and gives a nonzero
perverse eigenobject on the ordinary-flag stack.  Its singular
support lies in the required cone by Theorem~\ref{det:field}.
No regularity of any local nilpotent is used.
\end{proof}
